\documentclass[11pt, a4paper]{article}
% ── Shared preamble for all RTOS lab documents ────────────────────────────────
% Usage: % ── Shared preamble for all RTOS lab documents ────────────────────────────────
% Usage: % ── Shared preamble for all RTOS lab documents ────────────────────────────────
% Usage: \input{../preamble}  (from each labNN/ directory)
% ──────────────────────────────────────────────────────────────────────────────

% ── Packages ──────────────────────────────────────────────────────────────────
\usepackage[a4paper, top=2.5cm, bottom=2.5cm, left=2.5cm, right=2.5cm, headheight=28pt]{geometry}
\usepackage[utf8]{inputenc}
\usepackage[T1]{fontenc}
\usepackage{lmodern}
\usepackage{booktabs}
\usepackage{array}
\usepackage{tabularx}
\usepackage{longtable}
\usepackage{multirow}
\usepackage{colortbl}
\usepackage{xcolor}
\usepackage{titlesec}
\usepackage{enumitem}
\usepackage{hyperref}
\usepackage{fancyhdr}
\usepackage{parskip}
\usepackage{listings}
\usepackage[most]{tcolorbox}
\usepackage{fontawesome5}
\usepackage{microtype}
\usepackage{graphicx}
\usepackage{amsmath}

% ── Colors (KMUTNB brand) ──────────────────────────────────────────────────────
\definecolor{kmutnbBlue}{RGB}{0,56,116}
\definecolor{kmutnbGold}{RGB}{186,146,36}
\definecolor{lightGray}{RGB}{245,245,245}
\definecolor{midGray}{RGB}{200,200,200}
\definecolor{codeGray}{RGB}{248,248,248}
\definecolor{codeFrame}{RGB}{220,220,220}
\definecolor{tipteal}{RGB}{0,105,92}
\definecolor{warnAmber}{RGB}{121,85,0}
\definecolor{checkGreen}{RGB}{27,94,32}

% ── Section formatting ─────────────────────────────────────────────────────────
\titleformat{\section}
  {\large\bfseries\color{kmutnbBlue}}
  {\thesection.}{0.5em}{}
  [\color{kmutnbGold}\titlerule]

\titleformat{\subsection}
  {\normalsize\bfseries\color{kmutnbBlue}}
  {\thesubsection.}{0.5em}{}

\titleformat{\subsubsection}
  {\small\bfseries\color{kmutnbBlue}}
  {\thesubsubsection.}{0.5em}{}

\titlespacing*{\section}{0pt}{1.4ex plus .2ex}{0.6ex plus .1ex}
\titlespacing*{\subsection}{0pt}{1.0ex plus .2ex}{0.4ex plus .1ex}

% ── Listings (C and bash) ──────────────────────────────────────────────────────
\lstdefinestyle{cstyle}{
  language        = C,
  basicstyle      = \ttfamily\small,
  keywordstyle    = \bfseries\color{kmutnbBlue},
  commentstyle    = \itshape\color{gray},
  stringstyle     = \color{tipteal},
  numberstyle     = \tiny\color{midGray},
  numbers         = left,
  numbersep       = 6pt,
  stepnumber      = 1,
  backgroundcolor = \color{codeGray},
  frame           = single,
  framerule       = 0.4pt,
  rulecolor       = \color{codeFrame},
  breaklines      = true,
  showstringspaces= false,
  tabsize         = 4,
  xleftmargin     = 1.5em,
  xrightmargin    = 0.5em,
  aboveskip       = 0.8em,
  belowskip       = 0.8em,
}

\lstdefinestyle{bashstyle}{
  language        = bash,
  basicstyle      = \ttfamily\small,
  keywordstyle    = \color{kmutnbBlue},
  commentstyle    = \itshape\color{gray},
  backgroundcolor = \color{black!5},
  frame           = single,
  framerule       = 0.4pt,
  rulecolor       = \color{codeFrame},
  breaklines      = true,
  showstringspaces= false,
  xleftmargin     = 1.5em,
  xrightmargin    = 0.5em,
  aboveskip       = 0.6em,
  belowskip       = 0.6em,
}

\lstdefinestyle{textstyle}{
  basicstyle      = \ttfamily\footnotesize,
  backgroundcolor = \color{black!4},
  frame           = single,
  framerule       = 0.4pt,
  rulecolor       = \color{codeFrame},
  breaklines      = true,
  xleftmargin     = 1.5em,
  xrightmargin    = 0.5em,
  aboveskip       = 0.6em,
  belowskip       = 0.6em,
}

% ── Callout boxes ──────────────────────────────────────────────────────────────
\tcbuselibrary{skins, breakable}

\newtcolorbox{tipbox}[1][]{
  enhanced, breakable,
  colback=tipteal!8, colframe=tipteal, coltitle=white,
  fonttitle=\bfseries\small, title={\faLightbulb\quad Tip#1},
  left=6pt, right=6pt, top=4pt, bottom=4pt,
  boxrule=0.8pt, arc=3pt,
}

\newtcolorbox{warnbox}[1][]{
  enhanced, breakable,
  colback=warnAmber!8, colframe=kmutnbGold, coltitle=white,
  fonttitle=\bfseries\small, title={\faExclamationTriangle\quad Note#1},
  left=6pt, right=6pt, top=4pt, bottom=4pt,
  boxrule=0.8pt, arc=3pt,
}

\newtcolorbox{checkpointbox}[1][]{
  enhanced,
  colback=kmutnbBlue!6, colframe=kmutnbBlue, coltitle=white,
  fonttitle=\bfseries\small, title={\faCheckCircle\quad Checkpoint#1},
  left=6pt, right=6pt, top=4pt, bottom=4pt,
  boxrule=0.8pt, arc=3pt,
}

\newtcolorbox{questionbox}[1][]{
  enhanced, breakable,
  colback=kmutnbGold!10, colframe=kmutnbGold, coltitle=white,
  fonttitle=\bfseries\small, title={\faQuestion\quad Question#1},
  left=6pt, right=6pt, top=4pt, bottom=4pt,
  boxrule=0.8pt, arc=3pt,
}

% ── Checkbox list ──────────────────────────────────────────────────────────────
\newlist{checklist}{itemize}{1}
\setlist[checklist]{label=\faSquare[regular], leftmargin=1.5em, noitemsep}

% ── Misc helpers ───────────────────────────────────────────────────────────────
\newcommand{\file}[1]{\texttt{\small #1}}
\newcommand{\api}[1]{\texttt{\small #1}}

% ── Title block macro ─────────────────────────────────────────────────────────
% Usage: \labtitle{03}{Aperiodic Servers \& Producer--Consumer}{2}{CLO 1 \& CLO 2}
\newcommand{\labtitle}[4]{%
  \begin{center}
    {\color{kmutnbBlue}\rule{\linewidth}{2pt}}\\[6pt]
    {\LARGE\bfseries\color{kmutnbBlue} Laboratory #1}\\[4pt]
    {\large\bfseries #2}\\[2pt]
    {\normalsize Real-Time Operating Systems \enspace\textbullet\enspace
     M.Eng.\ ECE \enspace\textbullet\enspace KMUTNB}\\[8pt]
    {\color{kmutnbGold}\rule{\linewidth}{1pt}}
  \end{center}
  \vspace{0.3cm}
  \begin{tabular}{@{}p{3.8cm} p{10cm}@{}}
    \textbf{Platform}       & NXP FRDM-MCXN236 (ARM Cortex-M33 @ 150\,MHz) \\[2pt]
    \textbf{RTOS}           & FreeRTOS v10.6 \\[2pt]
    \textbf{Toolchain}      & ARM GNU Toolchain (\texttt{arm-none-eabi-gcc}) + Make \\[2pt]
    \textbf{Estimated time} & #3 hours \\[2pt]
    \textbf{CLO mapping}    & #4 \\
  \end{tabular}
  \vspace{0.4cm}
}
  (from each labNN/ directory)
% ──────────────────────────────────────────────────────────────────────────────

% ── Packages ──────────────────────────────────────────────────────────────────
\usepackage[a4paper, top=2.5cm, bottom=2.5cm, left=2.5cm, right=2.5cm, headheight=28pt]{geometry}
\usepackage[utf8]{inputenc}
\usepackage[T1]{fontenc}
\usepackage{lmodern}
\usepackage{booktabs}
\usepackage{array}
\usepackage{tabularx}
\usepackage{longtable}
\usepackage{multirow}
\usepackage{colortbl}
\usepackage{xcolor}
\usepackage{titlesec}
\usepackage{enumitem}
\usepackage{hyperref}
\usepackage{fancyhdr}
\usepackage{parskip}
\usepackage{listings}
\usepackage[most]{tcolorbox}
\usepackage{fontawesome5}
\usepackage{microtype}
\usepackage{graphicx}
\usepackage{amsmath}

% ── Colors (KMUTNB brand) ──────────────────────────────────────────────────────
\definecolor{kmutnbBlue}{RGB}{0,56,116}
\definecolor{kmutnbGold}{RGB}{186,146,36}
\definecolor{lightGray}{RGB}{245,245,245}
\definecolor{midGray}{RGB}{200,200,200}
\definecolor{codeGray}{RGB}{248,248,248}
\definecolor{codeFrame}{RGB}{220,220,220}
\definecolor{tipteal}{RGB}{0,105,92}
\definecolor{warnAmber}{RGB}{121,85,0}
\definecolor{checkGreen}{RGB}{27,94,32}

% ── Section formatting ─────────────────────────────────────────────────────────
\titleformat{\section}
  {\large\bfseries\color{kmutnbBlue}}
  {\thesection.}{0.5em}{}
  [\color{kmutnbGold}\titlerule]

\titleformat{\subsection}
  {\normalsize\bfseries\color{kmutnbBlue}}
  {\thesubsection.}{0.5em}{}

\titleformat{\subsubsection}
  {\small\bfseries\color{kmutnbBlue}}
  {\thesubsubsection.}{0.5em}{}

\titlespacing*{\section}{0pt}{1.4ex plus .2ex}{0.6ex plus .1ex}
\titlespacing*{\subsection}{0pt}{1.0ex plus .2ex}{0.4ex plus .1ex}

% ── Listings (C and bash) ──────────────────────────────────────────────────────
\lstdefinestyle{cstyle}{
  language        = C,
  basicstyle      = \ttfamily\small,
  keywordstyle    = \bfseries\color{kmutnbBlue},
  commentstyle    = \itshape\color{gray},
  stringstyle     = \color{tipteal},
  numberstyle     = \tiny\color{midGray},
  numbers         = left,
  numbersep       = 6pt,
  stepnumber      = 1,
  backgroundcolor = \color{codeGray},
  frame           = single,
  framerule       = 0.4pt,
  rulecolor       = \color{codeFrame},
  breaklines      = true,
  showstringspaces= false,
  tabsize         = 4,
  xleftmargin     = 1.5em,
  xrightmargin    = 0.5em,
  aboveskip       = 0.8em,
  belowskip       = 0.8em,
}

\lstdefinestyle{bashstyle}{
  language        = bash,
  basicstyle      = \ttfamily\small,
  keywordstyle    = \color{kmutnbBlue},
  commentstyle    = \itshape\color{gray},
  backgroundcolor = \color{black!5},
  frame           = single,
  framerule       = 0.4pt,
  rulecolor       = \color{codeFrame},
  breaklines      = true,
  showstringspaces= false,
  xleftmargin     = 1.5em,
  xrightmargin    = 0.5em,
  aboveskip       = 0.6em,
  belowskip       = 0.6em,
}

\lstdefinestyle{textstyle}{
  basicstyle      = \ttfamily\footnotesize,
  backgroundcolor = \color{black!4},
  frame           = single,
  framerule       = 0.4pt,
  rulecolor       = \color{codeFrame},
  breaklines      = true,
  xleftmargin     = 1.5em,
  xrightmargin    = 0.5em,
  aboveskip       = 0.6em,
  belowskip       = 0.6em,
}

% ── Callout boxes ──────────────────────────────────────────────────────────────
\tcbuselibrary{skins, breakable}

\newtcolorbox{tipbox}[1][]{
  enhanced, breakable,
  colback=tipteal!8, colframe=tipteal, coltitle=white,
  fonttitle=\bfseries\small, title={\faLightbulb\quad Tip#1},
  left=6pt, right=6pt, top=4pt, bottom=4pt,
  boxrule=0.8pt, arc=3pt,
}

\newtcolorbox{warnbox}[1][]{
  enhanced, breakable,
  colback=warnAmber!8, colframe=kmutnbGold, coltitle=white,
  fonttitle=\bfseries\small, title={\faExclamationTriangle\quad Note#1},
  left=6pt, right=6pt, top=4pt, bottom=4pt,
  boxrule=0.8pt, arc=3pt,
}

\newtcolorbox{checkpointbox}[1][]{
  enhanced,
  colback=kmutnbBlue!6, colframe=kmutnbBlue, coltitle=white,
  fonttitle=\bfseries\small, title={\faCheckCircle\quad Checkpoint#1},
  left=6pt, right=6pt, top=4pt, bottom=4pt,
  boxrule=0.8pt, arc=3pt,
}

\newtcolorbox{questionbox}[1][]{
  enhanced, breakable,
  colback=kmutnbGold!10, colframe=kmutnbGold, coltitle=white,
  fonttitle=\bfseries\small, title={\faQuestion\quad Question#1},
  left=6pt, right=6pt, top=4pt, bottom=4pt,
  boxrule=0.8pt, arc=3pt,
}

% ── Checkbox list ──────────────────────────────────────────────────────────────
\newlist{checklist}{itemize}{1}
\setlist[checklist]{label=\faSquare[regular], leftmargin=1.5em, noitemsep}

% ── Misc helpers ───────────────────────────────────────────────────────────────
\newcommand{\file}[1]{\texttt{\small #1}}
\newcommand{\api}[1]{\texttt{\small #1}}

% ── Title block macro ─────────────────────────────────────────────────────────
% Usage: \labtitle{03}{Aperiodic Servers \& Producer--Consumer}{2}{CLO 1 \& CLO 2}
\newcommand{\labtitle}[4]{%
  \begin{center}
    {\color{kmutnbBlue}\rule{\linewidth}{2pt}}\\[6pt]
    {\LARGE\bfseries\color{kmutnbBlue} Laboratory #1}\\[4pt]
    {\large\bfseries #2}\\[2pt]
    {\normalsize Real-Time Operating Systems \enspace\textbullet\enspace
     M.Eng.\ ECE \enspace\textbullet\enspace KMUTNB}\\[8pt]
    {\color{kmutnbGold}\rule{\linewidth}{1pt}}
  \end{center}
  \vspace{0.3cm}
  \begin{tabular}{@{}p{3.8cm} p{10cm}@{}}
    \textbf{Platform}       & NXP FRDM-MCXN236 (ARM Cortex-M33 @ 150\,MHz) \\[2pt]
    \textbf{RTOS}           & FreeRTOS v10.6 \\[2pt]
    \textbf{Toolchain}      & ARM GNU Toolchain (\texttt{arm-none-eabi-gcc}) + Make \\[2pt]
    \textbf{Estimated time} & #3 hours \\[2pt]
    \textbf{CLO mapping}    & #4 \\
  \end{tabular}
  \vspace{0.4cm}
}
  (from each labNN/ directory)
% ──────────────────────────────────────────────────────────────────────────────

% ── Packages ──────────────────────────────────────────────────────────────────
\usepackage[a4paper, top=2.5cm, bottom=2.5cm, left=2.5cm, right=2.5cm, headheight=28pt]{geometry}
\usepackage[utf8]{inputenc}
\usepackage[T1]{fontenc}
\usepackage{lmodern}
\usepackage{booktabs}
\usepackage{array}
\usepackage{tabularx}
\usepackage{longtable}
\usepackage{multirow}
\usepackage{colortbl}
\usepackage{xcolor}
\usepackage{titlesec}
\usepackage{enumitem}
\usepackage{hyperref}
\usepackage{fancyhdr}
\usepackage{parskip}
\usepackage{listings}
\usepackage[most]{tcolorbox}
\usepackage{fontawesome5}
\usepackage{microtype}
\usepackage{graphicx}
\usepackage{amsmath}

% ── Colors (KMUTNB brand) ──────────────────────────────────────────────────────
\definecolor{kmutnbBlue}{RGB}{0,56,116}
\definecolor{kmutnbGold}{RGB}{186,146,36}
\definecolor{lightGray}{RGB}{245,245,245}
\definecolor{midGray}{RGB}{200,200,200}
\definecolor{codeGray}{RGB}{248,248,248}
\definecolor{codeFrame}{RGB}{220,220,220}
\definecolor{tipteal}{RGB}{0,105,92}
\definecolor{warnAmber}{RGB}{121,85,0}
\definecolor{checkGreen}{RGB}{27,94,32}

% ── Section formatting ─────────────────────────────────────────────────────────
\titleformat{\section}
  {\large\bfseries\color{kmutnbBlue}}
  {\thesection.}{0.5em}{}
  [\color{kmutnbGold}\titlerule]

\titleformat{\subsection}
  {\normalsize\bfseries\color{kmutnbBlue}}
  {\thesubsection.}{0.5em}{}

\titleformat{\subsubsection}
  {\small\bfseries\color{kmutnbBlue}}
  {\thesubsubsection.}{0.5em}{}

\titlespacing*{\section}{0pt}{1.4ex plus .2ex}{0.6ex plus .1ex}
\titlespacing*{\subsection}{0pt}{1.0ex plus .2ex}{0.4ex plus .1ex}

% ── Listings (C and bash) ──────────────────────────────────────────────────────
\lstdefinestyle{cstyle}{
  language        = C,
  basicstyle      = \ttfamily\small,
  keywordstyle    = \bfseries\color{kmutnbBlue},
  commentstyle    = \itshape\color{gray},
  stringstyle     = \color{tipteal},
  numberstyle     = \tiny\color{midGray},
  numbers         = left,
  numbersep       = 6pt,
  stepnumber      = 1,
  backgroundcolor = \color{codeGray},
  frame           = single,
  framerule       = 0.4pt,
  rulecolor       = \color{codeFrame},
  breaklines      = true,
  showstringspaces= false,
  tabsize         = 4,
  xleftmargin     = 1.5em,
  xrightmargin    = 0.5em,
  aboveskip       = 0.8em,
  belowskip       = 0.8em,
}

\lstdefinestyle{bashstyle}{
  language        = bash,
  basicstyle      = \ttfamily\small,
  keywordstyle    = \color{kmutnbBlue},
  commentstyle    = \itshape\color{gray},
  backgroundcolor = \color{black!5},
  frame           = single,
  framerule       = 0.4pt,
  rulecolor       = \color{codeFrame},
  breaklines      = true,
  showstringspaces= false,
  xleftmargin     = 1.5em,
  xrightmargin    = 0.5em,
  aboveskip       = 0.6em,
  belowskip       = 0.6em,
}

\lstdefinestyle{textstyle}{
  basicstyle      = \ttfamily\footnotesize,
  backgroundcolor = \color{black!4},
  frame           = single,
  framerule       = 0.4pt,
  rulecolor       = \color{codeFrame},
  breaklines      = true,
  xleftmargin     = 1.5em,
  xrightmargin    = 0.5em,
  aboveskip       = 0.6em,
  belowskip       = 0.6em,
}

% ── Callout boxes ──────────────────────────────────────────────────────────────
\tcbuselibrary{skins, breakable}

\newtcolorbox{tipbox}[1][]{
  enhanced, breakable,
  colback=tipteal!8, colframe=tipteal, coltitle=white,
  fonttitle=\bfseries\small, title={\faLightbulb\quad Tip#1},
  left=6pt, right=6pt, top=4pt, bottom=4pt,
  boxrule=0.8pt, arc=3pt,
}

\newtcolorbox{warnbox}[1][]{
  enhanced, breakable,
  colback=warnAmber!8, colframe=kmutnbGold, coltitle=white,
  fonttitle=\bfseries\small, title={\faExclamationTriangle\quad Note#1},
  left=6pt, right=6pt, top=4pt, bottom=4pt,
  boxrule=0.8pt, arc=3pt,
}

\newtcolorbox{checkpointbox}[1][]{
  enhanced,
  colback=kmutnbBlue!6, colframe=kmutnbBlue, coltitle=white,
  fonttitle=\bfseries\small, title={\faCheckCircle\quad Checkpoint#1},
  left=6pt, right=6pt, top=4pt, bottom=4pt,
  boxrule=0.8pt, arc=3pt,
}

\newtcolorbox{questionbox}[1][]{
  enhanced, breakable,
  colback=kmutnbGold!10, colframe=kmutnbGold, coltitle=white,
  fonttitle=\bfseries\small, title={\faQuestion\quad Question#1},
  left=6pt, right=6pt, top=4pt, bottom=4pt,
  boxrule=0.8pt, arc=3pt,
}

% ── Checkbox list ──────────────────────────────────────────────────────────────
\newlist{checklist}{itemize}{1}
\setlist[checklist]{label=\faSquare[regular], leftmargin=1.5em, noitemsep}

% ── Misc helpers ───────────────────────────────────────────────────────────────
\newcommand{\file}[1]{\texttt{\small #1}}
\newcommand{\api}[1]{\texttt{\small #1}}

% ── Title block macro ─────────────────────────────────────────────────────────
% Usage: \labtitle{03}{Aperiodic Servers \& Producer--Consumer}{2}{CLO 1 \& CLO 2}
\newcommand{\labtitle}[4]{%
  \begin{center}
    {\color{kmutnbBlue}\rule{\linewidth}{2pt}}\\[6pt]
    {\LARGE\bfseries\color{kmutnbBlue} Laboratory #1}\\[4pt]
    {\large\bfseries #2}\\[2pt]
    {\normalsize Real-Time Operating Systems \enspace\textbullet\enspace
     M.Eng.\ ECE \enspace\textbullet\enspace KMUTNB}\\[8pt]
    {\color{kmutnbGold}\rule{\linewidth}{1pt}}
  \end{center}
  \vspace{0.3cm}
  \begin{tabular}{@{}p{3.8cm} p{10cm}@{}}
    \textbf{Platform}       & NXP FRDM-MCXN236 (ARM Cortex-M33 @ 150\,MHz) \\[2pt]
    \textbf{RTOS}           & FreeRTOS v10.6 \\[2pt]
    \textbf{Toolchain}      & ARM GNU Toolchain (\texttt{arm-none-eabi-gcc}) + Make \\[2pt]
    \textbf{Estimated time} & #3 hours \\[2pt]
    \textbf{CLO mapping}    & #4 \\
  \end{tabular}
  \vspace{0.4cm}
}


% ── Header / Footer ────────────────────────────────────────────────────────────
\pagestyle{fancy}
\fancyhf{}
\fancyhead[L]{\small\color{kmutnbBlue}\textbf{KMUTNB -- Faculty of Engineering}}
\fancyhead[R]{\small\color{kmutnbBlue}Dept.\ of Electrical and Computer Engineering}
\fancyfoot[L]{\small\color{kmutnbBlue}Real-Time Operating Systems -- M.Eng.\ ECE}
\fancyfoot[C]{\small\thepage}
\fancyfoot[R]{\small\color{kmutnbBlue}Lab 08}
\renewcommand{\headrulewidth}{0.4pt}
\renewcommand{\headrule}{\hbox to\headwidth{%
  \color{kmutnbGold}\leaders\hrule height\headrulewidth\hfill}}
\renewcommand{\footrulewidth}{0.4pt}
\renewcommand{\footrule}{\hbox to\headwidth{%
  \color{midGray}\leaders\hrule height\footrulewidth\hfill}}

\hypersetup{
  colorlinks = true,
  linkcolor  = kmutnbBlue,
  urlcolor   = kmutnbBlue,
  citecolor  = kmutnbBlue,
  pdftitle   = {Lab 08 -- TrustZone-M: Secure/Non-Secure Partition},
  pdfauthor  = {KMUTNB RTOS Course},
}

% ══════════════════════════════════════════════════════════════════════════════
\begin{document}

\labtitle{08}{TrustZone-M: Secure/Non-Secure Partition}{4--5}{CLO 5 \& CLO 6}

% ── Objectives ────────────────────────────────────────────────────────────────
\section{Objectives}

By completing this lab you will be able to:

\begin{enumerate}[noitemsep, leftmargin=*]
  \item \textbf{Configure} the SAU to partition flash and SRAM into Secure, NSC, and
        Non-Secure regions.
  \item \textbf{Implement} a Secure API using the CMSE
        \texttt{\_\_attribute\_\_((cmse\_nonsecure\_entry))} annotation.
  \item \textbf{Call} a Secure function from a Non-Secure FreeRTOS task via an NSC veneer.
  \item \textbf{Trigger} a SecureFault by attempting a direct NS$\to$Secure SRAM access.
  \item \textbf{Measure} NSC call overhead vs.\ a direct function call using DWT.
\end{enumerate}

% ── Prerequisites ─────────────────────────────────────────────────────────────
\section{Prerequisites}

\begin{checklist}
  \item Lab 07 complete (stack analysis and static allocation working).
  \item ARM GNU Toolchain with CMSE support (\texttt{arm-none-eabi-gcc} $\geq$ 10,
        compiled with \texttt{-mcmse}).
  \item MCUXpresso IDE or command-line build with separate Secure and Non-Secure projects.
  \item SEGGER J-Link or MCU-Link (J-Link firmware) for debugging SecureFault.
\end{checklist}

% ── Background ────────────────────────────────────────────────────────────────
\section{Background}

\subsection{TrustZone-M Execution Model}

The Cortex-M33 has two security states: \textbf{Secure (S)} and \textbf{Non-Secure (NS)}.
The processor is always in one state. Transitions happen only through specific instructions:

\begin{itemize}[noitemsep, leftmargin=*]
  \item \textbf{NS $\to$ S:} only through the SG (Secure Gateway) instruction at an address
        in a Non-Secure Callable (NSC) region.
  \item \textbf{S $\to$ NS:} via \texttt{BXNS} / \texttt{BLXNS} instructions.
\end{itemize}

\subsection{Two-Project Build Structure}

TrustZone-M requires two separate firmware images:

\begin{lstlisting}[style=textstyle]
Secure project (linked to 0x0000_0000)
  - Runs first after reset
  - Configures SAU, TrustZone peripherals
  - Implements secure API functions
  - Releases Non-Secure image

Non-Secure project (linked to 0x0004_0000)
  - FreeRTOS scheduler
  - Application tasks
  - Calls Secure API through NSC veneers
\end{lstlisting}

% ── Project Structure ─────────────────────────────────────────────────────────
\section{Project Structure}

\begin{lstlisting}[style=textstyle]
lab08_trustzone/
  secure/                        <- Secure-world project
    Makefile
    linker_secure.ld             <- Secure flash + SRAM regions
    src/
      main_secure.c              <- SAU config, Secure init, NS release
      secure_api.c/.h            <- NSC entry functions
      secure_api_veneer.h        <- auto-generated veneer declarations
      startup_MCXN236_s.s        <- Secure vector table
    cmse_implib.o                <- output: NSC import library
  nonsecure/                     <- Non-Secure FreeRTOS project
    Makefile
    linker_nonsecure.ld          <- NS flash + SRAM regions
    src/
      main_ns.c                  <- FreeRTOS init, task creation
      ns_tasks.c/.h              <- tasks that call Secure API
      FreeRTOSConfig.h
      uart.h / uart.c
    FreeRTOS-Kernel/
  flash_both.sh                  <- flashes S image then NS image
  src/
    uart.h                       <- shared UART declaration
\end{lstlisting}

% ── Part A ─────────────────────────────────────────────────────────────────────
\section{Part A --- SAU Configuration}

\subsection{Step 1 --- Define the Memory Partition}

Edit \file{secure/src/main\_secure.c} to configure the Security Attribution Unit:

\begin{lstlisting}[style=cstyle]
void SAU_Configure(void)
{
    /* Disable SAU before configuration */
    SAU->CTRL = 0;

    /* Region 0: Non-Secure flash [0x0004_0000 .. 0x001F_FFFF] */
    SAU->RNR  = 0;
    SAU->RBAR = 0x00040000U;
    SAU->RLAR = (0x001FFFFFU & SAU_RLAR_LADDR_Msk) | SAU_RLAR_ENABLE_Msk;

    /* Region 1: NSC region in flash [0x0003_FE00 .. 0x0003_FFFF] (512 B) */
    SAU->RNR  = 1;
    SAU->RBAR = 0x0003FE00U;
    SAU->RLAR = (0x0003FFFFU & SAU_RLAR_LADDR_Msk)
              | SAU_RLAR_NSC_Msk | SAU_RLAR_ENABLE_Msk;

    /* Region 2: Non-Secure SRAM [0x2002_0000 .. 0x200F_FFFF] */
    SAU->RNR  = 2;
    SAU->RBAR = 0x20020000U;
    SAU->RLAR = (0x200FFFFFU & SAU_RLAR_LADDR_Msk) | SAU_RLAR_ENABLE_Msk;

    /* Enable SAU */
    SAU->CTRL = SAU_CTRL_ENABLE_Msk;
    __DSB();
    __ISB();
}
\end{lstlisting}

\subsection{Step 2 --- Verify with \texttt{TZM\_IsSecure}}

Add a diagnostic function that probes whether an address is Secure:

\begin{lstlisting}[style=cstyle]
/* Secure-side only -- NS code cannot call this */
bool TZM_IsSecure(uint32_t addr)
{
    for (int i = 0; i < 8; i++) {
        SAU->RNR = i;
        if (!(SAU->RLAR & SAU_RLAR_ENABLE_Msk)) continue;
        uint32_t base  = SAU->RBAR & SAU_RBAR_BADDR_Msk;
        uint32_t limit = (SAU->RLAR & SAU_RLAR_LADDR_Msk) | 0x1F;
        if (addr >= base && addr <= limit) return false;
    }
    return true; /* not in any NS/NSC region -> Secure */
}
\end{lstlisting}

Call from \file{main\_secure.c} before releasing the NS core:

\begin{lstlisting}[style=cstyle]
uart_printf("[S] 0x00010000 is %s\r\n",
    TZM_IsSecure(0x00010000) ? "SECURE" : "NS/NSC");
uart_printf("[S] 0x00050000 is %s\r\n",
    TZM_IsSecure(0x00050000) ? "SECURE" : "NS/NSC");
uart_printf("[S] 0x0003FE00 is %s (NSC)\r\n",
    TZM_IsSecure(0x0003FE00) ? "SECURE" : "NS/NSC");
\end{lstlisting}

\begin{checkpointbox}[~A]
  UART output showing correct Secure/NS attribution for the three test addresses.
\end{checkpointbox}

% ── Part B ─────────────────────────────────────────────────────────────────────
\section{Part B --- Implement and Call Secure API}

\subsection{Step 1 --- Write the Secure API}

In \file{secure/src/secure\_api.c}:

\begin{lstlisting}[style=cstyle]
#include <arm_cmse.h>
#include "secure_api.h"

/* Secure ADC value -- NS world must not read this directly */
static volatile int32_t s_last_adc_value = 0;

/* NSC entry function -- callable from Non-Secure world */
__attribute__((cmse_nonsecure_entry))
int32_t Secure_ADC_Read(void)
{
    s_last_adc_value = (s_last_adc_value + 37) % 4096;
    return s_last_adc_value;
}

__attribute__((cmse_nonsecure_entry))
void Secure_ADC_Init(void)
{
    uart_printf("[S] Secure_ADC_Init called from NS world\r\n");
    s_last_adc_value = 0;
}
\end{lstlisting}

In \file{secure/src/secure\_api.h}:

\begin{lstlisting}[style=cstyle]
#ifndef SECURE_API_H
#define SECURE_API_H
#include <stdint.h>

/* Declarations for NS world -- placed in CMSE import library */
int32_t Secure_ADC_Read(void);
void    Secure_ADC_Init(void);

#endif /* SECURE_API_H */
\end{lstlisting}

\subsection{Step 2 --- Build the Secure Image and Generate the Import Library}

\begin{lstlisting}[style=bashstyle]
cd secure
make
# Output: secure.elf  +  cmse_implib.o (NSC veneer import library)
\end{lstlisting}

The linker flags \texttt{-Wl,\allowbreak--cmse-implib -Wl,\allowbreak--out-implib=cmse\_implib.o}
generate the import library containing SG stubs.

Verify SG stubs appear in the map file:

\begin{lstlisting}[style=bashstyle]
grep ".gnu.sgstubs" secure.map
\end{lstlisting}

Expected output:
\begin{lstlisting}[style=textstyle]
.gnu.sgstubs   0x0003fe00   0x10   Secure_ADC_Read
.gnu.sgstubs   0x0003fe10   0x08   Secure_ADC_Init
\end{lstlisting}

\begin{checkpointbox}[~B1]
  Map file excerpt showing SG stubs located in the NSC region (0x0003FE00).
\end{checkpointbox}

\subsection{Step 3 --- Call from Non-Secure FreeRTOS Task}

Copy \file{cmse\_implib.o} to \file{nonsecure/}. In \file{nonsecure/src/ns\_tasks.c}:

\begin{lstlisting}[style=cstyle]
#include "secure_api.h"

void vSensorTask(void *pv)
{
    Secure_ADC_Init();
    for (;;) {
        int32_t value = Secure_ADC_Read();
        uart_printf("[NS] ADC value = %ld\r\n", value);
        vTaskDelay(pdMS_TO_TICKS(200));
    }
}
\end{lstlisting}

Build the NS image (links against \file{cmse\_implib.o}):

\begin{lstlisting}[style=bashstyle]
cd nonsecure
make
\end{lstlisting}

\subsection{Step 4 --- Flash and Run}

\begin{lstlisting}[style=bashstyle]
./flash_both.sh
# Flashes secure.bin to 0x0000_0000, then nonsecure.bin to 0x0004_0000
\end{lstlisting}

Expected UART output:

\begin{lstlisting}[style=textstyle]
[S] SAU configured
[S] 0x00010000 is SECURE
[S] 0x00050000 is NS/NSC
[S] Releasing Non-Secure image at 0x00040000
[NS] starting FreeRTOS...
[S] Secure_ADC_Init called from NS world
[NS] ADC value = 37
[NS] ADC value = 74
\end{lstlisting}

\begin{checkpointbox}[~B2]
  UART output showing alternating \texttt{[S]} init and \texttt{[NS]} sensor values.
\end{checkpointbox}

% ── Part C ─────────────────────────────────────────────────────────────────────
\section{Part C --- Trigger a SecureFault}

\subsection{Step 1 --- Direct NS$\to$Secure SRAM Access}

In \texttt{vSensorTask}, add a deliberate access to Secure SRAM (this will fault):

\begin{lstlisting}[style=cstyle]
/* Attempting to read Secure SRAM from NS world -- should SecureFault */
volatile uint32_t *secure_ptr = (volatile uint32_t *)0x20010000U;
uint32_t val = *secure_ptr;   /* <- triggers SecureFault */
uart_printf("[NS] should never print: %lu\r\n", val);
\end{lstlisting}

\subsection{Step 2 --- Implement the SecureFault Handler}

In \file{nonsecure/src/main\_ns.c}:

\begin{lstlisting}[style=cstyle]
void SecureFault_Handler(void)
{
    uart_printf("[NS] !!! SecureFault !!!\r\n");
    uart_printf("[NS] SFSR = 0x%08lX\r\n", SAU->SFSR);
    uart_printf("[NS] SFAR = 0x%08lX\r\n", SAU->SFAR);
    for (;;) {}
}
\end{lstlisting}

\begin{checkpointbox}[~C]
  UART showing \texttt{!!! SecureFault !!!} with SFSR/SFAR values.
  Note the faulting address (0x20010000) reported in SFAR.
\end{checkpointbox}

\begin{questionbox}[~C1]
  The SecureFault is raised on the \textbf{NS} side (not the Secure side). Explain
  why --- which hardware unit detects the violation and in which security state does
  the fault handler run?
\end{questionbox}

\begin{questionbox}[~C2]
  The SFSR (Secure Fault Status Register) has several bits. Look up the meaning of
  bit 1 (\texttt{SFARVALID}) and bit 0 (\texttt{INVEP}). Which bit is set in your
  fault? What does it indicate?
\end{questionbox}

% ── Part D ─────────────────────────────────────────────────────────────────────
\section{Part D --- Measure NSC Call Overhead}

\subsection{Step 1 --- Benchmark NSC vs.\ Direct Call}

Use the DWT cycle counter (Lab 06) to compare NSC and direct call latency:

\begin{lstlisting}[style=cstyle]
void vBenchmarkTask(void *pv)
{
    DWT_Enable();
    uint32_t nsc_cycles, direct_cycles;

    /* NSC call overhead */
    DWT_START();
    volatile int32_t v1 = Secure_ADC_Read();   /* crosses S/NS boundary */
    DWT_STOP();
    nsc_cycles = _dwt_elapsed;

    /* Direct function call (NS-only function, same logic) */
    DWT_START();
    volatile int32_t v2 = ns_dummy_adc_read();
    DWT_STOP();
    direct_cycles = _dwt_elapsed;

    uart_printf("[BENCH] NSC: %lu cycles  Direct: %lu cycles  Overhead: %lu\r\n",
                nsc_cycles, direct_cycles, nsc_cycles - direct_cycles);
    vTaskDelete(NULL);
}
\end{lstlisting}

\begin{checkpointbox}[~D]
  Terminal output showing NSC and direct call cycle counts with computed overhead.
\end{checkpointbox}

\begin{questionbox}[~D1]
  The NSC call overhead includes \texttt{BLXNS} + SG instruction + context save.
  At 150\,MHz, convert the overhead cycles to nanoseconds. Is this acceptable for
  a 1\,kHz control loop?
\end{questionbox}

\begin{questionbox}[~D2]
  FreeRTOS-M33 (NS kernel) allocates a ``Secure context'' per NS task that calls NSC
  functions. Where is this context stored, and why is it needed?
\end{questionbox}

% ── Deliverables ───────────────────────────────────────────────────────────────
\section{Deliverables}

\renewcommand{\arraystretch}{1.3}
\begin{tabular}{@{} c p{4.5cm} p{7cm} @{}}
  \toprule
  \textbf{\#} & \textbf{Item} & \textbf{Details} \\
  \midrule
  1 & Part A UART & SAU probe output --- 3 addresses \\
  2 & Part B1 map excerpt & SG stubs in NSC region \\
  3 & Part B2 UART & Alternating \texttt{[S]} init and \texttt{[NS]} sensor values \\
  4 & Part C UART & SecureFault with SFSR/SFAR values \\
  5 & Part D UART & NSC vs.\ direct call cycle counts \\
  6 & Written answers & Questions C1--C2, D1--D2 \\
  7 & Source listings & \file{secure\_api.c} and \file{ns\_tasks.c} \\
  \bottomrule
\end{tabular}
\renewcommand{\arraystretch}{1}

% ── Key Concepts ───────────────────────────────────────────────────────────────
\section{Key Concepts Demonstrated}

\renewcommand{\arraystretch}{1.3}
\begin{tabular}{@{} p{7cm} p{6cm} @{}}
  \toprule
  \textbf{Concept} & \textbf{Where you saw it} \\
  \midrule
  SAU region configuration & Part A \\
  Security attribution verification & Part A \texttt{TZM\_IsSecure} \\
  CMSE \texttt{cmse\_nonsecure\_entry} + SG veneer & Part B \\
  NSC import library and BLXNS call & Part B \\
  SecureFault on illegal NS$\to$S access & Part C \\
  SFSR / SFAR fault diagnosis & Part C \\
  NSC call overhead measurement & Part D \\
  \bottomrule
\end{tabular}
\renewcommand{\arraystretch}{1}

% ── Reference ─────────────────────────────────────────────────────────────────
\section{Reference}

\renewcommand{\arraystretch}{1.3}
\begin{tabular}{@{} p{5.5cm} p{7.5cm} @{}}
  \toprule
  \textbf{Resource} & \textbf{Relevant sections} \\
  \midrule
  ARM Cortex-M33 TRM & \S B9 TrustZone-M, \S B2.3 SecureFault \\
  NXP AN13814 & TrustZone-M on MCXN236 \\
  ARM CMSE specification & \texttt{cmse\_nonsecure\_entry}, \texttt{BLXNS} \\
  GCC ARM options & \texttt{-mcmse}, \texttt{--cmse-implib}, \texttt{--out-implib} \\
  \bottomrule
\end{tabular}
\renewcommand{\arraystretch}{1}

\end{document}
